% GENERATED FILE. Edit canonical lessons, then run npm run generate:books.
% canonical-source-hash: fnv1a64:991a2274bd3837ae
% canonical-lessons: ZH-C321-lishi, ZH-C321-guren, ZH-W321-gu-ancient, ZH-C321-gushu, ZH-C321-gulao, ZH-C321-wenzi

\chapter{People of Long Ago}
\label{ch:zh-people-of-long-ago}
\hlchaptermodality{321}

\begin{chapteropening}
\textbf{By the end of this chapter:} \emph{I can talk about history, the people of long ago, old books and old scripts.}

Complete the last lesson of chapter 321: I can talk about history, the people of long ago, old books and old scripts.
\end{chapteropening}

\section[lìshǐ]{lìshǐ — history}

\label{lesson:ZH-C321-lishi}



\begin{quote}
Before the new one: say the Mandarin for to bring up a subject, then the Mandarin for a voice message.
\end{quote}

\subsection*{You'll want to know: lìshǐ}
\emph{lìshǐ} — \textquotedblleft{}history\textquotedblright{}.

\begin{quote}
\emph{wǒ zài dàxué xué guo Zhōngguó de lìshǐ} — I studied Chinese history at university.
\end{quote}

Its characters come later; for now it is a word for the ear.

\begin{cousinweb}
\emph{lìshǐ} is built from \emph{lì} 'to pass through' and \emph{shǐ} 'a record', both inherited from Middle Chinese, \emph{lek} and \emph{sriX} (Guangyun, in Baxter's transcription), which glosses the first 'to go through' and calls the second 'one who records events'.
\end{cousinweb}

\begin{grammarlens}[title={What you've built}]
One more word for telling what happened.
\end{grammarlens}

\subsection*{Your turn}
\begin{itemize}
\raggedright
  \item \emph{Say it:} \emph{lìshǐ}
  \item \emph{Say it:} \emph{lìshǐ}, once more
  \item \emph{Say it:} \emph{wǒ zài dàxué xué guo Zhōngguó de lìshǐ} — I studied Chinese history at university.
  \item \emph{Recall:} say the Mandarin for to bring up a subject, then the Mandarin for a voice message, then say \emph{lìshǐ} again
\end{itemize}

\begin{tcolorbox}[breakable,colback=teal!4,colframe=teal!35!black,title={Before you move on}]
What does lìshǐ mean? (\textquotedblleft{}History\textquotedblright{}.) Say it once more.
\end{tcolorbox}
\section[\texorpdfstring{\zh{古人} (gǔrén)}{古人 (gǔrén)}]{\zh{古人} (gǔrén) — the ancients; people of long ago}

\label{lesson:ZH-C321-guren}



\begin{quote}
Before the new one: say the Mandarin for a voice message, then the Mandarin for the history of a country.
\end{quote}

\subsection*{You'll want to know: \zh{古人}}
\textbf{\zh{古人}} — \emph{gǔrén} — \textquotedblleft{}the ancients; people of long ago\textquotedblright{}.

\begin{quote}
\textbf{\zh{古人没有电话}} \emph{gǔrén méiyǒu diànhuà} — People long ago had no telephones.
\end{quote}

The new character, in pieces: \zh{古} — five strokes: \zh{十} on top, and \zh{口} under it. The next lesson puts it in your hand.

\begin{cousinweb}
\textbf{\zh{古人}} (\emph{gǔrén}) is built from \zh{古} 'ancient' and \zh{人} 'a person', both inherited from Middle Chinese, \emph{kuX} and \emph{nyin} (Guangyun, in Baxter's transcription), which glosses the first 'former'.
\end{cousinweb}

\begin{grammarlens}[title={What you've built}]
One more word for telling what happened.
\end{grammarlens}

\subsection*{Your turn}
\begin{itemize}
\raggedright
  \item \emph{Say it:} \emph{gǔrén}
  \item \emph{Say it:} \emph{gǔrén}, once more
  \item \emph{Say it:} \emph{gǔrén méiyǒu diànhuà} — People long ago had no telephones.
  \item \emph{Recall:} say the Mandarin for a voice message, then the Mandarin for the history of a country, then say \emph{gǔrén} again
\end{itemize}

\begin{tcolorbox}[breakable,colback=teal!4,colframe=teal!35!black,title={Before you move on}]
What does \zh{古人} mean? (\textquotedblleft{}The ancients; people of long ago\textquotedblright{}.) Say it once more.
\end{tcolorbox}
\section[\texorpdfstring{\zh{古} (gǔ)}{古 (gǔ)}]{\zh{古} — \zh{十} over \zh{口}}

\label{lesson:ZH-W321-gu-ancient}



\begin{quote}
The word first, then the character inside it.

\begin{itemize}
\raggedright
  \item \emph{Recall:} say \emph{gǔrén} — \textbf{R1}, one lesson back, before it decays
\end{itemize}

\begin{quote}
\zh{十} \zh{口}
\end{quote}

\emph{Read it:} each one aloud, with nothing said first — \emph{shí}, \emph{kǒu}
\end{quote}

\subsection*{Script — \zh{十} over \zh{口}}
\hlblockfigure{\detokenize{figures/ZH-W321-gu-ancient-filmstrip.pdf}}{How \zh{古} is written, stroke by stroke}

\begin{quote}
\zh{古}
\end{quote}

\textbf{Two pieces, top to bottom.}

\textbf{The top piece, \zh{十}, two strokes.} A long bar from left to right, then a vertical down through its middle.

\textbf{The bottom piece, \zh{口}, three strokes.} A short vertical on the left; a bar that runs right and turns down; then the bar that closes the foot.

2 + 3 = five.

\subsection*{Your turn}
\begin{itemize}
\raggedright
  \item \emph{Write it:} \zh{古}, all five, with the model in front of you
  \item \emph{Cover:} write it again from memory, and count the pieces, not the strokes
\end{itemize}

\begin{tcolorbox}[breakable,colback=teal!4,colframe=teal!35!black,title={Before you move on}]
How many strokes in \zh{古}? (\textbf{Five}.) What is on top? (\textbf{\zh{十}}.) What is under it? (\textbf{\zh{口}}.)

Source: \href{https://raw.githubusercontent.com/chanind/hanzi-writer-data/68d10a4b21150cae5e1ebbd223eed289cf32d90c/data/\%E5\%8F\%A4.json}{Hanzi Writer Data \zh{古}, pinned stroke paths and medians}.
\end{tcolorbox}
\section[\texorpdfstring{\zh{古书} (gǔshū)}{古书 (gǔshū)}]{\zh{古书} (gǔshū) — ancient book; old book}

\label{lesson:ZH-C321-gushu}



\begin{quote}
Before the new one: say the Mandarin for the history of a country, then the Mandarin for the people of long ago.
\end{quote}

\subsection*{You'll want to know: \zh{古书}}
\textbf{\zh{古书}} — \emph{gǔshū} — \textquotedblleft{}ancient book; old book\textquotedblright{}.

\begin{quote}
\textbf{\zh{我在书店里看到了一本古书}} \emph{wǒ zài shūdiàn lǐ kàndào le yì běn gǔshū} — I came across an old book in the bookshop.
\end{quote}

\begin{cousinweb}
\textbf{\zh{古书}} (\emph{gǔshū}) is built from \zh{古} 'ancient' and \zh{书} 'a book', both inherited from Middle Chinese, \emph{kuX} and \emph{syo} (Guangyun, in Baxter's transcription), which glosses the first 'former' and says of the second that it sets things down.
\end{cousinweb}

\begin{grammarlens}[title={What you've built}]
One more word for telling what happened.
\end{grammarlens}

\subsection*{Your turn}
\begin{itemize}
\raggedright
  \item \emph{Say it:} \emph{gǔshū}
  \item \emph{Say it:} \emph{gǔshū}, once more
  \item \emph{Say it:} \emph{wǒ zài shūdiàn lǐ kàndào le yì běn gǔshū} — I came across an old book in the bookshop.
  \item \emph{Recall:} say the Mandarin for the history of a country, then the Mandarin for the people of long ago, then say \emph{gǔshū} again
\end{itemize}

\begin{tcolorbox}[breakable,colback=teal!4,colframe=teal!35!black,title={Before you move on}]
What does \zh{古书} mean? (\textquotedblleft{}Ancient book; old book\textquotedblright{}.) Say it once more.
\end{tcolorbox}
\section[\texorpdfstring{\zh{古老} (gǔlǎo)}{古老 (gǔlǎo)}]{\zh{古老} (gǔlǎo) — ancient; age-old}

\label{lesson:ZH-C321-gulao}



\begin{quote}
Before the new one: say the Mandarin for the people of long ago, then the Mandarin for an old book from long ago.
\end{quote}

\subsection*{You'll want to know: \zh{古老}}
\textbf{\zh{古老}} — \emph{gǔlǎo} — \textquotedblleft{}ancient; age-old\textquotedblright{}.

\begin{quote}
\textbf{\zh{他们家的房子很古老}} \emph{tāmen jiā de fángzi hěn gǔlǎo} — Their family's house is ancient.
\end{quote}

Two third tones meet in this word: gǔ is said with a rising tone before lǎo, and the spelling keeps both marks.

\begin{cousinweb}
\textbf{\zh{古老}} (\emph{gǔlǎo}) is built from \zh{古} 'ancient' and \zh{老} 'old', both inherited from Middle Chinese, \emph{kuX} and \emph{lawX} (Guangyun, in Baxter's transcription), which glosses the first 'former' and the second 'aged'.
\end{cousinweb}

\begin{grammarlens}[title={What you've built}]
One more word for telling what happened.
\end{grammarlens}

\subsection*{Your turn}
\begin{itemize}
\raggedright
  \item \emph{Say it:} \emph{gǔlǎo}
  \item \emph{Say it:} \emph{gǔlǎo}, once more
  \item \emph{Say it:} \emph{tāmen jiā de fángzi hěn gǔlǎo} — Their family's house is ancient.
  \item \emph{Recall:} say the Mandarin for the people of long ago, then the Mandarin for an old book from long ago, then say \emph{gǔlǎo} again
\end{itemize}

\begin{tcolorbox}[breakable,colback=teal!4,colframe=teal!35!black,title={Before you move on}]
What does \zh{古老} mean? (\textquotedblleft{}Ancient; age-old\textquotedblright{}.) Say it once more.
\end{tcolorbox}
\section[\texorpdfstring{\zh{文字} (wénzì)}{文字 (wénzì)}]{\zh{文字} (wénzì) — script; written language}

\label{lesson:ZH-C321-wenzi}



\begin{quote}
Before the new one: say the Mandarin for an old book from long ago, then the Mandarin for age-old.
\end{quote}

\subsection*{You'll want to know: \zh{文字}}
\textbf{\zh{文字}} — \emph{wénzì} — \textquotedblleft{}script; written language\textquotedblright{}.

\begin{quote}
\textbf{\zh{汉字是很古老的文字}} \emph{Hànzì shì hěn gǔlǎo de wénzì} — Chinese characters are a very old script.
\end{quote}

\begin{cousinweb}
\textbf{\zh{文字}} (\emph{wénzì}) is built from \zh{文} 'writing' and \zh{字} 'a character', both inherited from Middle Chinese, \emph{mjun} and \emph{dziH} (Guangyun, in Baxter's transcription), which glosses the first 'writing' and 'beautiful'; for the second it gives 'an ornament' and 'to nurse', and the sense 'a written character' is not given there.
\end{cousinweb}

\begin{grammarlens}[title={What you've built}]
One more word for telling what happened.
\end{grammarlens}

\subsection*{Your turn}
\begin{itemize}
\raggedright
  \item \emph{Say it:} \emph{wénzì}
  \item \emph{Say it:} \emph{wénzì}, once more
  \item \emph{Say it:} \emph{Hànzì shì hěn gǔlǎo de wénzì} — Chinese characters are a very old script.
  \item \emph{Recall:} say the Mandarin for an old book from long ago, then the Mandarin for age-old, then say \emph{wénzì} again
\end{itemize}

\begin{tcolorbox}[breakable,colback=teal!4,colframe=teal!35!black,title={Before you move on}]
What does \zh{文字} mean? (\textquotedblleft{}Script; written language\textquotedblright{}.) Say it once more.
\end{tcolorbox}
